% ============================================================
% Chapitre 3 : Estimation Ponctuelle — EMV, Methode des Moments
% ============================================================
\chapter{Estimation Ponctuelle --- EMV et M\'ethode des Moments}
\label{ch:estimation}

\begin{center}
\textit{<<~L'art de la statistique est de choisir le bon r\'esum\'e num\'erique.~>>}
\end{center}

% ------------------------------------------------------------
\section{Exemple motivant}
% ------------------------------------------------------------

On mesure la dur\'ee de vie (en heures) de 10 ampoules LED :
\[
1200,\; 1350,\; 980,\; 1500,\; 1100,\; 1450,\; 1300,\; 1050,\; 1400,\; 1250.
\]
On mod\'elise par une loi exponentielle $\mathcal{E}(\lambda)$. Comment estimer $\lambda$ ? Deux m\'ethodes classiques s'offrent \`a nous : la \textbf{m\'ethode des moments} et le \textbf{maximum de vraisemblance}.

% ------------------------------------------------------------
\section{Cadre g\'en\'eral}
% ------------------------------------------------------------

\begin{definition}[Estimateur]
\index{estimateur}
Soit $(X_1,\ldots,X_n)\iid P_\theta$. Un \textbf{estimateur} de $\theta$ (ou d'une fonction $g(\theta)$) est une statistique $\thetahat = T(X_1,\ldots,X_n)$ \`a valeurs dans $\Theta$.
\end{definition}

\begin{definition}[Estimation]
L'\textbf{estimation} $\hat{\theta}=T(x_1,\ldots,x_n)$ est la valeur num\'erique obtenue en rempla\c{c}ant les v.a.\ par les observations.
\end{definition}

% ------------------------------------------------------------
\section{M\'ethode des moments}
\label{sec:moments}
% ------------------------------------------------------------

\begin{definition}[M\'ethode des moments (MoM)]
\index{methode des moments@m\'ethode des moments}
On \'egale les $k$ premiers moments th\'eoriques aux moments empiriques :
\[
\E_\theta[X^j] = \frac{1}{n}\sum_{i=1}^n X_i^j, \quad j=1,\ldots,k,
\]
o\`u $k=\dim(\theta)$, puis on r\'esout le syst\`eme en $\theta$.
\end{definition}

\begin{exemple}[Loi normale]
Soit $X_1,\ldots,X_n\iid\mathcal{N}(\mu,\sigma^2)$. On a $\E[X]=\mu$ et $\E[X^2]=\mu^2+\sigma^2$. La MoM donne :
\[
\hat{\mu}_{\mathrm{MoM}} = \xbar, \qquad \hat{\sigma}^2_{\mathrm{MoM}} = \frac{1}{n}\sum_{i=1}^n X_i^2 - \xbar^2 = \frac{1}{n}\sum_{i=1}^n (X_i-\xbar)^2.
\]
\end{exemple}

\begin{exemple}[Loi exponentielle]
$X\sim\mathcal{E}(\lambda)$ : $\E[X]=1/\lambda$. Donc $\hat{\lambda}_{\mathrm{MoM}} = 1/\xbar$.
\end{exemple}

\begin{exemple}[Loi Gamma]
$X\sim\Gamma(\alpha,\beta)$ avec $\E[X]=\alpha/\beta$ et $\Var(X)=\alpha/\beta^2$. On obtient :
\[
\hat{\beta}_{\mathrm{MoM}} = \frac{\xbar}{S^2}, \qquad \hat{\alpha}_{\mathrm{MoM}} = \frac{\xbar^2}{S^2}.
\]
\end{exemple}

\begin{intuition}
La MoM est simple et donne souvent de bons estimateurs, mais elle n'utilise pas toute l'information contenue dans l'\'echantillon. Le maximum de vraisemblance est g\'en\'eralement pr\'ef\'erable.
\end{intuition}

% ------------------------------------------------------------
\section{Maximum de vraisemblance}
\label{sec:emv}
% ------------------------------------------------------------

\begin{definition}[Fonction de vraisemblance]
\index{vraisemblance}
La \textbf{vraisemblance} de l'\'echantillon $(x_1,\ldots,x_n)$ est :
\[
L(\theta) = L(\theta;x_1,\ldots,x_n) = \prod_{i=1}^n f(x_i;\theta),
\]
o\`u $f(\cdot;\theta)$ est la densit\'e (cas continu) ou la fonction de masse (cas discret).

La \textbf{log-vraisemblance} est $\ell(\theta)=\ln L(\theta) = \sum_{i=1}^n \ln f(x_i;\theta)$.
\end{definition}

\begin{definition}[Estimateur du maximum de vraisemblance (EMV)]
\index{EMV}\index{maximum de vraisemblance}
L'\textbf{EMV} est :
\[
\thetahat_{\mathrm{EMV}} = \argmax_{\theta\in\Theta} L(\theta) = \argmax_{\theta\in\Theta} \ell(\theta).
\]
\end{definition}

\begin{remarque}
En pratique, on r\'esout souvent l'\textbf{\'equation de vraisemblance} :
\[
\frac{\partial \ell}{\partial \theta}(\theta) = 0,
\]
en v\'erifiant que la solution est bien un maximum (condition d'ordre 2 ou argument de convexit\'e).
\end{remarque}

\subsection{Exemples fondamentaux}

\begin{exemple}[Loi de Bernoulli]
$X_1,\ldots,X_n\iid\mathcal{B}(1,p)$. On a $\ell(p)=\sum x_i\ln p + (n-\sum x_i)\ln(1-p)$. L'\'equation $\ell'(p)=0$ donne :
\[
\frac{\sum x_i}{p} - \frac{n-\sum x_i}{1-p} = 0 \implies \phat_{\mathrm{EMV}} = \frac{1}{n}\sum_{i=1}^n X_i = \xbar.
\]
\end{exemple}

\begin{exemple}[Loi normale]
\label{ex:emv-normale}
$X_1,\ldots,X_n\iid\mathcal{N}(\mu,\sigma^2)$.
\[
\ell(\mu,\sigma^2) = -\frac{n}{2}\ln(2\pi) - \frac{n}{2}\ln\sigma^2 - \frac{1}{2\sigma^2}\sum_{i=1}^n(x_i-\mu)^2.
\]
Les \'equations $\partial\ell/\partial\mu=0$ et $\partial\ell/\partial\sigma^2=0$ donnent :
\[
\hat{\mu}_{\mathrm{EMV}} = \xbar, \qquad \hat{\sigma}^2_{\mathrm{EMV}} = \frac{1}{n}\sum_{i=1}^n(X_i-\xbar)^2.
\]
\end{exemple}

\begin{attention}
L'EMV de $\sigma^2$ sous le mod\`ele normal est $\frac{1}{n}\sum(X_i-\xbar)^2$ (diviseur $n$), qui est \textbf{biais\'e}. La version corrig\'ee $S^2$ (diviseur $n-1$) est sans biais (cf.\ chapitre~\ref{ch:proprietes}).
\end{attention}

\begin{exemple}[Loi exponentielle]
$X_1,\ldots,X_n\iid\mathcal{E}(\lambda)$. On a $\ell(\lambda)=n\ln\lambda-\lambda\sum x_i$. L'\'equation $\ell'(\lambda)=0$ donne $\hat{\lambda}_{\mathrm{EMV}}=n/\sum X_i = 1/\xbar$.
\end{exemple}

\begin{exemple}[Loi de Poisson]
$X_1,\ldots,X_n\iid\mathcal{P}(\lambda)$. On a $\ell(\lambda)=\sum x_i\ln\lambda - n\lambda - \sum\ln(x_i!)$. L'\'equation $\ell'(\lambda)=0$ donne $\hat{\lambda}_{\mathrm{EMV}}=\xbar$.
\end{exemple}

\begin{exemple}[Loi uniforme]
$X_1,\ldots,X_n\iid\mathcal{U}([0,\theta])$. La vraisemblance est $L(\theta)=\theta^{-n}\mathbf{1}_{\theta\geq x_{(n)}}$, d\'ecroissante en $\theta$ pour $\theta\geq x_{(n)}$. Donc $\thetahat_{\mathrm{EMV}}=X_{(n)}=\max_i X_i$.

Ici, l'\'equation de vraisemblance ne s'applique pas (le maximum est atteint au bord).
\end{exemple}

% ------------------------------------------------------------
\section{Propri\'et\'es de l'EMV}
% ------------------------------------------------------------

\begin{theoreme}[Propri\'et\'e d'invariance]
\index{invariance de l'EMV}
Si $\thetahat$ est l'EMV de $\theta$, alors pour toute fonction $g$, l'EMV de $g(\theta)$ est $g(\thetahat)$.
\end{theoreme}

\begin{theoreme}[Consistance de l'EMV]
Sous des conditions de r\'egularit\'e (famille exponentielle, etc.), l'EMV est \textbf{convergent} :
\[
\thetahat_{\mathrm{EMV}} \xrightarrow{\Prob} \theta_0 \quad (n\to\infty).
\]
\end{theoreme}

\begin{theoreme}[Normalit\'e asymptotique de l'EMV]
\label{thm:normalite-emv}
Sous des conditions de r\'egularit\'e :
\[
\sqrt{n}(\thetahat_{\mathrm{EMV}}-\theta_0) \xrightarrow{\mathcal{L}} \mathcal{N}\!\left(0, \frac{1}{I(\theta_0)}\right),
\]
o\`u $I(\theta)=-\E\!\left[\frac{\partial^2 \ln f(X;\theta)}{\partial\theta^2}\right]$ est l'\textbf{information de Fisher} (cf.\ chapitre~\ref{ch:proprietes}).
\end{theoreme}

% ------------------------------------------------------------
\section{Information de Fisher}
% ------------------------------------------------------------

\begin{definition}[Information de Fisher]
\index{information de Fisher}
Pour un mod\`ele r\'egulier, l'\textbf{information de Fisher} d'une observation est :
\[
I(\theta) = \E_\theta\!\left[\left(\frac{\partial \ln f(X;\theta)}{\partial\theta}\right)^{\!2}\right] = -\E_\theta\!\left[\frac{\partial^2 \ln f(X;\theta)}{\partial\theta^2}\right].
\]
Pour un \'echantillon de taille $n$ : $I_n(\theta)=nI(\theta)$.
\end{definition}

\begin{exemple}
$X\sim\mathcal{N}(\mu,\sigma^2)$ avec $\sigma^2$ connu. $\ln f = -\frac{(x-\mu)^2}{2\sigma^2}+\text{cst}$, donc $\frac{\partial^2\ln f}{\partial\mu^2}=-\frac{1}{\sigma^2}$ et $I(\mu)=\frac{1}{\sigma^2}$.
\end{exemple}

\begin{exemple}
$X\sim\mathcal{E}(\lambda)$. $\ln f = \ln\lambda - \lambda x$, $\frac{\partial^2\ln f}{\partial\lambda^2}=-\frac{1}{\lambda^2}$, donc $I(\lambda)=\frac{1}{\lambda^2}$.
\end{exemple}

% ------------------------------------------------------------
\section{Impl\'ementation Python}
% ------------------------------------------------------------

\begin{python}
\begin{minted}{python}
import numpy as np
from scipy import stats
from scipy.optimize import minimize_scalar
import matplotlib.pyplot as plt

# --- Donnees : duree de vie d'ampoules ---
data = np.array([1200, 1350, 980, 1500, 1100, 1450, 1300, 1050, 1400, 1250])

# --- Methode des moments : Exp(lambda) ---
lambda_mom = 1 / np.mean(data)
print(f"MoM: lambda_hat = {lambda_mom:.6f}")

# --- EMV : Exp(lambda) ---
lambda_mle = 1 / np.mean(data)  # coincide avec MoM pour l'exponentielle
print(f"EMV: lambda_hat = {lambda_mle:.6f}")

# --- EMV pour la loi normale ---
mu_mle = np.mean(data)
sigma2_mle = np.var(data, ddof=0)  # diviseur n (EMV)
sigma2_unbiased = np.var(data, ddof=1)  # diviseur n-1 (sans biais)
print(f"EMV Normale: mu = {mu_mle:.2f}, sigma^2 = {sigma2_mle:.2f}")
print(f"Sans biais:  sigma^2 = {sigma2_unbiased:.2f}")

# --- EMV avec scipy (methode generale) ---
# Ajustement d'une loi exponentielle
loc, scale = stats.expon.fit(data, floc=0)  # floc=0 fixe le decalage
lambda_scipy = 1 / scale
print(f"scipy EMV: lambda = {lambda_scipy:.6f}")

# --- Visualisation de la log-vraisemblance ---
lambdas = np.linspace(0.0005, 0.002, 200)
log_lik = np.array([np.sum(stats.expon.logpdf(data, scale=1/lam)) for lam in lambdas])

plt.figure(figsize=(8, 5))
plt.plot(lambdas, log_lik, 'b-', lw=2)
plt.axvline(lambda_mle, color='r', linestyle='--', label=f'EMV = {lambda_mle:.5f}')
plt.xlabel(r'$\lambda$')
plt.ylabel(r'$\ell(\lambda)$')
plt.title('Log-vraisemblance exponentielle')
plt.legend()
plt.savefig("ch03_loglik.pdf")
plt.show()

# --- EMV loi uniforme U([0, theta]) ---
data_unif = np.array([0.3, 0.7, 0.5, 0.9, 0.2, 0.8, 0.6, 0.4, 0.95, 0.1])
theta_mle = np.max(data_unif)
print(f"EMV U([0,theta]): theta_hat = {theta_mle}")

# --- Information de Fisher ---
# Pour Exp(lambda) : I(lambda) = 1/lambda^2
n = len(data)
I_fisher = 1 / lambda_mle**2
var_asymp = 1 / (n * I_fisher)
print(f"Variance asymptotique de l'EMV: {var_asymp:.8f}")
print(f"IC asymptotique 95%: [{lambda_mle - 1.96*np.sqrt(var_asymp):.5f}, "
      f"{lambda_mle + 1.96*np.sqrt(var_asymp):.5f}]")
\end{minted}
\end{python}

% ------------------------------------------------------------
\section{Exercices}
% ------------------------------------------------------------

\begin{exercice}[$\star$ -- Calcul d'EMV]
Calculer l'EMV pour chacun des mod\`eles suivants :
\begin{enumerate}
\item $X_1,\ldots,X_n\iid\mathcal{G}\'eom\'etrique(p)$ (loi g\'eom\'etrique).
\item $X_1,\ldots,X_n\iid\mathcal{U}([a,b])$ avec $a$ et $b$ inconnus.
\end{enumerate}
\end{exercice}

\begin{exercice}[$\star$ -- MoM]
Soit $X_1,\ldots,X_n\iid\mathrm{Beta}(\alpha,\beta)$. On rappelle que $\E[X]=\frac{\alpha}{\alpha+\beta}$ et $\Var(X)=\frac{\alpha\beta}{(\alpha+\beta)^2(\alpha+\beta+1)}$. Trouver les estimateurs des moments pour $\alpha$ et $\beta$.
\end{exercice}

\begin{exercice}[$\star\star$ -- Invariance]
$X_1,\ldots,X_n\iid\mathcal{N}(\mu,\sigma^2)$. En utilisant l'invariance de l'EMV, d\'eterminer l'EMV de :
\begin{enumerate}
\item $\sigma$ (l'\'ecart-type),
\item $\Prob(X>c)$ pour $c$ fix\'e,
\item le coefficient de variation $\sigma/\mu$.
\end{enumerate}
\end{exercice}

\begin{exercice}[$\star\star$ -- Information de Fisher]
\begin{enumerate}
\item Calculer $I(\theta)$ pour $X\sim\mathcal{B}(1,p)$.
\item Calculer $I(\theta)$ pour $X\sim\mathcal{P}(\lambda)$.
\item Calculer la matrice d'information de Fisher $I(\mu,\sigma^2)$ pour $X\sim\mathcal{N}(\mu,\sigma^2)$.
\end{enumerate}
\end{exercice}

\begin{exercice}[$\star\star$ -- EMV non r\'egulier]
Soit $X_1,\ldots,X_n\iid\mathcal{U}([0,\theta])$.
\begin{enumerate}
\item Montrer que l'EMV est $\thetahat=X_{(n)}$.
\item Trouver la loi de $X_{(n)}$ et calculer $\E[X_{(n)}]$. L'EMV est-il sans biais ?
\item Proposer un estimateur sans biais bas\'e sur $X_{(n)}$.
\end{enumerate}
\end{exercice}

\begin{exercice}[$\star\star\star$ -- Projet : comparaison MoM vs EMV]
Pour la loi Gamma$(\alpha,\beta)$ :
\begin{enumerate}
\item Calculer les estimateurs MoM et EMV (num\'eriquement pour l'EMV).
\item Par simulation (1000 r\'ep\'etitions, $n=20,50,200$), comparer le biais et l'EQM des deux m\'ethodes.
\item Tracer les distributions d'\'echantillonnage des estimateurs et conclure.
\end{enumerate}
\end{exercice}

% ------------------------------------------------------------
\begin{formules}
\begin{align*}
\text{MoM} &: \E_\theta[X^j] = \frac{1}{n}\sum X_i^j,\; j=1,\ldots,\dim(\theta) \\
L(\theta) &= \prod_{i=1}^n f(x_i;\theta), \quad \ell(\theta)=\sum_{i=1}^n \ln f(x_i;\theta) \\
\thetahat_{\mathrm{EMV}} &= \argmax_\theta \ell(\theta) \\
I(\theta) &= -\E\!\left[\frac{\partial^2\ln f}{\partial\theta^2}\right] = \E\!\left[\left(\frac{\partial\ln f}{\partial\theta}\right)^{\!2}\right] \\
\sqrt{n}(\thetahat_{\mathrm{EMV}}-\theta) &\xrightarrow{\mathcal{L}} \mathcal{N}(0,1/I(\theta))
\end{align*}
\end{formules}
